\documentclass[10pt,a4paper]{amsart}
\usepackage[margin=19mm]{geometry}
\usepackage{amsmath,amssymb,mathtools}
\usepackage[scr=rsfs,cal=euler]{mathalfa}
\usepackage{xfrac}
\usepackage[unicode,hidelinks,bookmarks=false]{hyperref}
\newcommand{\ie}{i.e.\ }
\newcommand{\cf}{cf.\ }
\newcommand{\resp}{respectively\ }
\newcommand{\Subsection}{\subsection}
\providecommand{\nobreakdash}{\nobreak\hbox{-}}
\setlength{\emergencystretch}{2em}
\multlinegap0pt
\usepackage{bm}
\usepackage{bbm}
\usepackage{dsfont}
\usepackage{xspace}
\usepackage{cancel}
\usepackage{mathtools}
\usepackage{amsthm}
\newtheorem{theorem}{Theorem}
\newtheorem{proposition}[theorem]{Proposition}
\title[The noncommutative topological factor theorem for rank-one p]{The noncommutative topological factor theorem for rank-one product lattices (Proposition 2.1)}
\author{Cyril Houdayer, Corentin Le Bars}
\date{Source: arXiv:2608.07421v2 (CC BY 4.0)}
\begin{document}
\maketitle

\noindent\emph{Source and licence.} The noncommutative topological factor theorem for rank-one product lattices. Version arXiv:2608.07421v2 (CC BY 4.0), distributed under the Creative Commons Attribution 4.0 International licence, as recorded in the arXiv licence metadata for this exact version and shown on the abstract page https://arxiv.org/abs/2608.07421v2. Full rights notice: licenses/arxiv-2608.07421-CC-BY-4.0.txt.\par\smallskip

\noindent\textit{Editorial scope.} Counted unit: the Proposition printed as Proposition 2.1 of the article (source0.tex lines 447--471, source label \texttt{lem:abstract-recurrence-contraction}), delivered together with the article's own complete original proof of it (source0.tex lines 473--523, one \texttt{proof} environment of 51 lines). No mathematics is written, repaired or re-derived: statement and proof are the article's own text line for line. Required context retained uncounted: none. Every definition, display and label of those lines is retained unchanged. Omitted from this package and not redistributed: 1--446, 472--472, 524--1507. Nothing inside a retained excerpt is omitted or altered: every retained body is delivered exactly as the article's lines are written, so no omission receipt of any kind accompanies this package. Citations: the retained excerpts carry no citation command at all, so no bibliography entry of the article is transcribed and no reference is redistributed. Reference adaptation: none was needed and none was made; every reference inside the delivered unit resolves inside it, so no unresolved reference or citation is produced. Printed slips kept literal: no printed word, symbol, hypothesis or number of the counted statement or of its proof was altered. Layout and class: the article's own class is \texttt{amsart} with packages including amsmath, amsthm, amsfonts, amssymb, mathtools, enumerate, hyperref, graphicx, caption, subcaption, pgf, tikz, appendix, color; the wrapper is the lane's pinned \texttt{amsart} editorial wrapper at 10pt on A4, which restates the theorem-like environments the retained text needs and adds the article's own macro definitions that the retained text needs (none), with extra wrapper packages (bm, bbm, dsfont, xspace, cancel, mathtools). The delivered numbering is the unit's own. Assets: no retained span contains a figure environment, a graphics-inclusion command, a PDF-page inclusion, a table or a TikZ picture; no image file is copied into this package, the package declares no asset dependency and no image file is redistributed with it. Licence: version arXiv:2608.07421v2 of this article is distributed under the Creative Commons Attribution 4.0 International licence, as recorded in the arXiv licence metadata for this exact version and shown on the abstract page https://arxiv.org/abs/2608.07421v2. A full rights notice is delivered at licenses/arxiv-2608.07421-CC-BY-4.0.txt. Characters: every retained excerpt is pure ASCII, so the delivered engine, XeLaTeX with its default Latin Modern text font, reports no missing character.
	\begin{proposition}
		\label{lem:abstract-recurrence-contraction}
		\label{prop:codim-one-strip}
		Let $\varnothing\neq I\subsetneq N$ and put $J=N\setminus I$.  Let
		$U\subset G_I$ be an identity neighborhood.  For every $j\in J$, let
		$L_j\subsetneq X_j$ be compact and let $V_j\subset X_j$ be nonempty and open.
		Then there exists $\gamma\in\Gamma$ such that
		\begin{equation}
			\label{eq:simultaneous-recurrence-contraction}
			p_I(\gamma)\in U
			\quad\text{and}\quad
			p_j(\gamma)L_j\subset V_j
			\quad\text{for every }j\in J.
		\end{equation}
		
		In particular, if $j\in N$, $U\subset G_{N\setminus\{j\}}$ is an identity
		neighborhood, and $O,V\subset X_j$ are nonempty open sets with
		$X_j\setminus O\neq\varnothing$, then there exists $\gamma\in\Gamma$ such that
		\begin{equation}
			\label{eq:codim-one-strip}
			p_{N\setminus\{j\}}(\gamma)\in U,
			\qquad
			p_j(\gamma)(X_j\setminus O)\subset V.
		\end{equation}
	\end{proposition}

	\begin{proof}
		Fix $j\in J$. Choose $\xi_j^-\in X_j\setminus L_j$ and
		$\xi_j^+\in V_j\setminus\{\xi_j^-\}$.
		The second choice is possible because $X_j$ has no isolated points.  By
		conjugating a north--south element, choose $a_j\in G_j$ whose repelling and
		attracting fixed points are respectively $\xi_j^-$ and $\xi_j^+$.  Choose open
		neighborhoods $O_j\ni\xi_j^-$ and $W_j\ni\xi_j^+$ such that
		$\overline{O_j}\cap L_j=\varnothing$ and $\overline{W_j}\subset V_j$.
		By north--south dynamics, there is $m_j\geq1$ such that
		\begin{equation}
			\label{eq:abstract-ns-j}
			a_j^m(X_j\setminus O_j)\subset W_j
			\qquad\text{for every }m\geq m_j.
		\end{equation}
		Continuity of the action and compactness of $L_j$ and $\overline{W_j}$ give an
		identity neighborhood $B_j\subset G_j$ such that
		\begin{equation}
			\label{eq:abstract-small-perturbations}
			B_jL_j\subset X_j\setminus O_j,
			\qquad
			B_j^{-1}\overline{W_j}\subset V_j.
		\end{equation}
		
		Choose an identity neighborhood $B_I\subset G_I$ satisfying
		$B_I^{-1}B_I\subset U$, and put
		\[
		B=B_I\times\prod_{j\in J}B_j\subset G,
		\qquad
		a=(e_I,(a_j)_{j\in J})\in G.
		\]
		Let $\pi:G\to G/\Gamma$ be the quotient map.  The finite $G$-invariant measure
		on $G/\Gamma$ has full support, so the nonempty open set $\pi(B)$ has positive
		measure.  Left translation by $a$ is invertible and measure preserving.  By
		Poincar\'e recurrence, there are arbitrarily large integers $m$ such that
		$a^m\pi(B)\cap\pi(B)\neq\varnothing$.
		Choose such an $m$ with $m\geq\max_{j\in J}m_j$.  There are $b_1,b_2\in B$ such
		that $a^mb_1\Gamma=b_2\Gamma$. Hence
		$\gamma=b_2^{-1}a^mb_1\in\Gamma$. Writing
		$b_\ell=(b_{\ell,I},(b_{\ell,j})_{j\in J})$, we obtain
		$p_I(\gamma)=b_{2,I}^{-1}b_{1,I}\in B_I^{-1}B_I\subset U$.
		For $j\in J$, equations \eqref{eq:abstract-ns-j} and
		\eqref{eq:abstract-small-perturbations} give
		\[p_j(\gamma)L_j
			=b_{2,j}^{-1}a_j^mb_{1,j}L_j
			\subset B_j^{-1}a_j^m(X_j\setminus O_j)
			\subset B_j^{-1}W_j
			\subset V_j.\]
		This proves \eqref{eq:simultaneous-recurrence-contraction}. The final assertion
		is obtained by taking $I=N\setminus\{j\}$ and $L_j=X_j\setminus O$. Here $I$ is
		nonempty because $n\geq2$.
	\end{proof}

\end{document}
